% Options for packages loaded elsewhere
\PassOptionsToPackage{unicode}{hyperref}
\PassOptionsToPackage{hyphens}{url}
%
\documentclass[
  11pt,
]{article}
\usepackage{lmodern}
\usepackage{amssymb,amsmath}
\usepackage{ifxetex,ifluatex}
\ifnum 0\ifxetex 1\fi\ifluatex 1\fi=0 % if pdftex
  \usepackage[T1]{fontenc}
  \usepackage[utf8]{inputenc}
  \usepackage{textcomp} % provide euro and other symbols
\else % if luatex or xetex
  \usepackage{unicode-math}
  \defaultfontfeatures{Scale=MatchLowercase}
  \defaultfontfeatures[\rmfamily]{Ligatures=TeX,Scale=1}
\fi
% Use upquote if available, for straight quotes in verbatim environments
\IfFileExists{upquote.sty}{\usepackage{upquote}}{}
\IfFileExists{microtype.sty}{% use microtype if available
  \usepackage[]{microtype}
  \UseMicrotypeSet[protrusion]{basicmath} % disable protrusion for tt fonts
}{}
\makeatletter
\@ifundefined{KOMAClassName}{% if non-KOMA class
  \IfFileExists{parskip.sty}{%
    \usepackage{parskip}
  }{% else
    \setlength{\parindent}{0pt}
    \setlength{\parskip}{6pt plus 2pt minus 1pt}}
}{% if KOMA class
  \KOMAoptions{parskip=half}}
\makeatother
\usepackage{xcolor}
\IfFileExists{xurl.sty}{\usepackage{xurl}}{} % add URL line breaks if available
\IfFileExists{bookmark.sty}{\usepackage{bookmark}}{\usepackage{hyperref}}
\hypersetup{
  hidelinks,
  pdfcreator={LaTeX via pandoc}}
\urlstyle{same} % disable monospaced font for URLs
\usepackage[margin=1in]{geometry}
\setlength{\emergencystretch}{3em} % prevent overfull lines
\providecommand{\tightlist}{%
  \setlength{\itemsep}{0pt}\setlength{\parskip}{0pt}}
\setcounter{secnumdepth}{-\maxdimen} % remove section numbering
\usepackage{mathrsfs}
\usepackage{esint}
\setmathfont{Asana Math}
\setlength{\emergencystretch}{3em}

\author{}
\date{23 September 2026}

\begin{document}

\hypertarget{quantitative-non-saturation-of-scalarweyl-functionals-near-negative-einstein-metrics}{%
\section{Quantitative non-saturation of scalar--Weyl functionals near
negative Einstein
metrics}\label{quantitative-non-saturation-of-scalarweyl-functionals-near-negative-einstein-metrics}}

23 September 2026

\hypertarget{main-theorem}{%
\subsection{1. Main theorem}\label{main-theorem}}

Let \((M,g)\) be a closed connected oriented smooth four-manifold. Fix a
Spin\(^c\) structure, with real first Chern class \(a\), whose
unperturbed Seiberg--Witten equations have a solution for every metric
conformal to \(g\). The structure can depend on \(g\); it is fixed while
the conformal weight varies. Write \[
V=\operatorname{Vol}_g(M),\qquad p=(a_g^+)^2>0,\qquad
r=\fint_M|\operatorname{Ric}+g|^2,\qquad
b_a=\frac{(-a^2)_+}{V}.
\] Ricci uses the tensor norm, two-forms the exterior-algebra norm, and
\(W^+\) the Hilbert--Schmidt norm as an endomorphism of \(\Lambda^+\).
Set \(w=\sqrt6|W^+|\). For a constant \(\beta>0\), define the
dimensionless, nonnegative functional deficit \[
\mathfrak D_\beta(g,a)
=\frac{\beta+1}{\beta}\int_M(s^2+\beta w^2)
-72\pi^2p.
\tag{1}
\] Here dimensionless means invariant under constant metric rescaling;
\(\mathfrak D_\beta/V\), like \(r\), is measured after fixing the
Einstein target to \(-g\).

\textbf{Theorem 1 (functional non-saturation).} For every finite \(L>2\)
there is \(c_L>0\) such that, simultaneously for \(2\le\beta\le L\), \[
\boxed{
(1+b_a)\left(r+\frac{\mathfrak D_\beta(g,a)}V\right)
\ge c_L(\beta-2)^2.
}
\tag{2}
\] There is no scalar-curvature band, Sobolev constant,
injectivity-radius, diameter, volume, or full-curvature hypothesis.
Actual conformal monopole solvability and the displayed square cost are
the inputs.

If the selected class additionally satisfies
\(a^2\ge T:=2\chi(M)+3\tau(M)\), there is \(c'_L>0\) with \[
\boxed{
r+\frac{\mathfrak D_\beta(g,a)}V\ge c'_L(\beta-2)^2.
}
\tag{3}
\] In particular, for an exact metric \(\operatorname{Ric}=-g\) and such
a class, \[
\int(s^2+6\beta|W^+|^2)
\ge\frac{72\pi^2\beta}{\beta+1}p
+c'_L\frac{\beta}{\beta+1}(\beta-2)^2V.
\tag{4}
\]

The endpoint \(\beta=2\) must remain: a negative Kähler--Einstein
surface with its canonical Spin\(^c\) structure has \(r=0\) and
\(\mathfrak D_2=0\). Thus (2) is a compatibility theorem, not a
prohibition of the Kähler--Einstein branch.

The proof supplies all auxiliary forms and their joint analytic budget
from (1). It extends the former calibration at \(\beta=8\) to a
parameter family of curvature functionals. It does not, by itself,
identify a new smooth topology with no Einstein metric.

\hypertarget{the-exact-deficit-decomposition}{%
\subsection{2. The exact deficit
decomposition}\label{the-exact-deficit-decomposition}}

Put \[
q=-s+w,\quad q_+=\max(q,0),\quad q_-=\max(-q,0),\quad
\Delta_+=\int q_+^2-72\pi^2p.
\] The signed monopole estimate gives \(\Delta_+\ge0\). The relevant
classical inequality follows from the rescaled equations; the positive
part is needed when scalar curvature changes sign. Completing a square
gives the exact identity \[
\boxed{
\mathfrak D_\beta
=\Delta_++\int q_-^2+\frac1\beta\int(s+\beta w)^2.
}
\tag{5}
\] Indeed
\((\beta+1)(s^2+\beta w^2)/\beta =(-s+w)^2+(s+\beta w)^2/\beta\).

Write \(e=s+4\), \(z=s+\beta w\), and \[
q_\beta=4\frac{\beta+1}{\beta},\qquad
\kappa=r+\mathfrak D_\beta/V.
\] Then \[
w-4/\beta=(z-e)/\beta,\qquad
q-q_\beta=z/\beta-(1+1/\beta)e,\qquad
\int e^2\le4Vr.
\tag{6}
\] Consequently, uniformly for \(2\le\beta\le L\), \[
\int\bigl[(q-q_\beta)^2+(w-4/\beta)^2\bigr]
\le C_L V\kappa,
\quad
\int(|\operatorname{Ric}|^2+s^2+|W^+|^2)
\le C_L V(1+\kappa).
\tag{7}
\] These bounds control the Weyl norm, not its eigenspaces. The latter
come from the actual monopole below.

\hypertarget{actual-bounded-spinors-with-no-scalar-band}{%
\subsection{3. Actual bounded spinors with no scalar
band}\label{actual-bounded-spinors-with-no-scalar-band}}

Set \[
Q=\max(q,3),\qquad
f_0=\max\{\sqrt Q,(-s)_+/4\}.
\] Choose smooth positive majorants \(f_j\to f_0\) uniformly on the
fixed compact manifold. Conformal covariance, applied to
\(\widehat g=f_j^{-2}g\), supplies solutions \[
D_A\Phi=0,\quad -iF_A^+=f_j\sigma(\Phi),\quad
\psi=2\sqrt2\sigma(\Phi),\quad
\rho=|\Phi|^2=|\psi|,\quad u=3\rho/2.
\tag{8}
\] We use \(|\sigma(\Phi)|^2=|\Phi|^4/8\). The spinor equation is
\(4\nabla_A^*\nabla_A\Phi+(s+f_j\rho)\Phi=0\). At a maximum of \(\rho\),
\(f_j\rho\le-s\), so \[
0\le\rho\le4,\qquad 0\le u\le6.
\tag{9}
\] Every approximation error below tends to zero for the fixed \(g,a\).
No uniform smoothing rate along a sequence and no bounds on derivatives
of \(f_j\) are required.

The cutoff costs only scalar error: \[
0\le Q^2-q_+^2\le9e^2,\qquad
0\le f_0^2-Q\le\frac9{16}e^2.
\tag{10}
\] For the second inequality, a nonzero difference forces
\((-s)_+>4\sqrt3>6\), whence \((-s)_+\le3|e|\). For the first, \(q<3\)
implies \(s>-3\), hence \(e>1\).

Define nonnegative defects \[
\begin{split}
D_0={}&2\int|d\rho|^2
+\int(4\rho|\nabla_A\Phi|^2-|\nabla\psi|^2)\\
&+2\int|d\psi|^2+2\int E_\lambda,\qquad
E_\lambda=\langle W^+\psi,\psi\rangle+w\rho^2/3\ge0.
\end{split}
\tag{11}
\] Nonnegativity uses the differential inequality for the parallel
spinor quadratic map and \(\lambda_{\min}(W^+)\ge-w/3\). The spinor and
self-dual-form Weitzenböck identities, subtracted after testing the
spinor against \(\rho\Phi\), give \[
\frac{27}{8}D_0=\int(qu^2-f_j u^3).
\tag{12}
\]

For completeness, the moment argument that retains the truncation cost
is as follows. Set \[
C=\int Qu^2,\quad Z=\int\sqrt Q\,u^3,\quad E_Q=\int Q^2.
\] Then \(C-Z\ge27D_0/8\ge0\), while the curvature equation gives \[
C_{\rm act}:=\int f_j^2u^2=18\|F_A^+\|_2^2\ge72\pi^2p,
\quad 0\le C_{\rm act}-C\le\frac{81}{4}\int e^2+o(1).
\] Hölder gives \(C\le E_Q^{1/3}Z^{2/3}\), hence \[
0\le E_Q-C\le S+o(1),\qquad
S:=\Delta_++\frac{117}{4}\int e^2\le117V\kappa.
\] If \(\theta=(Z/E_Q)^{1/3}\), the nonnegative Hölder slack implies
\(E_Q-C\ge2(C-Z)\). The associated amplitude estimate is
\(\int Q(u-\sqrt Q)^2\le E_Q-C\): with probability measure
\(Q^2d\mu/E_Q\), let \(h=u/\sqrt Q\), \(c=\mathbb E h^2\),
\(z_3=\mathbb E h^3\), and \(m=\mathbb E h\). Then
\(z_3\le c\le z_3^{2/3}\), so \(c\le1\); Cauchy--Schwarz gives
\(c^2\le m z_3\le mc\), and hence \(m\ge c\). Thus
\(\mathbb E(h-1)^2\le1-c\). We have proved \[
D_0\le\frac{4S}{27}+o(1),\qquad
\int Q(u-\sqrt Q)^2\le S+o(1).
\tag{13}
\]

\hypertarget{produce-the-full-calibrated-form-including-its-weyl-direction}{%
\subsection{4. Produce the full calibrated form, including its Weyl
direction}\label{produce-the-full-calibrated-form-including-its-weyl-direction}}

Because \(Q\ge3\), \(u\le6\), and \((Q-q)^2\le e^2\), (13) yields \[
\int(q-u^2)^2
\le2(1+2\sqrt3)^2S+2\int e^2+o(1)
\le C V\kappa+o(1).
\tag{14}
\] The middle inequality uses
\(|Q-u^2|\le(1+2\sqrt3)\sqrt Q\,|\sqrt Q-u|\). Together with (7), this
controls \(\int(u^2-q_\beta)^2\).

Now choose the normalization, rather than assuming it: \[
c_\beta^2=\frac{9\beta}{8(\beta+1)},\qquad
\omega=c_\beta\psi,\qquad t_\beta=1/\beta,\qquad
\alpha_\beta=\frac9{8(\beta+1)}.
\] It follows that \[
\|\omega\|_\infty\le3\sqrt2,\qquad
\int(|\omega|^2-2)^2\le C_L V\kappa+o(1).
\tag{15}
\] Define the smooth polynomial model \[
A_{\omega,t}=t\left(\frac{|\omega|^2}{3}I-\omega\otimes\omega\right).
\] The exact identity for arbitrary \(\psi\), including its zeros, is \[
\left|W^+-\alpha_\beta
       \left(\frac{\rho^2}{3}I-\psi\otimes\psi\right)\right|^2
=\frac16(w-2\alpha_\beta\rho^2)^2
+2\alpha_\beta E_\lambda.
\tag{16}
\] This identity uses only trace-freeness of \(W^+\), not an eigenline
selection. The tensor on its left is exactly \(W^+-A_{\omega,1/\beta}\).
Also \(2\alpha_\beta\rho^2=u^2/(\beta+1)\), and \[
w-\frac{u^2}{\beta+1}
=\left(w-\frac4\beta\right)
-\frac{u^2-q_\beta}{\beta+1}.
\] Equations (7), (13)--(16) give, for the same smooth solution, \[
d_\beta:=\fint\left(
 |\operatorname{Ric}+g|^2+|W^+-A_{\omega,1/\beta}|^2
+(|\omega|^2-2)^2\right)
\le C_L\kappa+o(1).
\tag{17}
\] In particular the small Weyl-direction error has been proved from the
actual gauge equations. It is not a hidden geometric hypothesis.

\hypertarget{supply-the-second-order-budget}{%
\subsection{5. Supply the second-order
budget}\label{supply-the-second-order-budget}}

For a bounded section \(v\) of a metric bundle with compatible
connection curvature \(\mathscr F\), covariant integration by parts on a
closed four-manifold gives \[
\|\nabla v\|_4^2\le4\|v\|_\infty\|\nabla^2v\|_2,
\] \[
\|\nabla^2v\|_2^2
\le2\|\nabla^*\nabla v\|_2^2
+32\|v\|_\infty^2\|\operatorname{Ric}\|_2^2
+33\|v\|_\infty^2\|\mathscr F\|_2^2.
\tag{18}
\] The curvature convention is
\(|\mathscr F|^2=\sum_{ij} |\mathscr F_{ij}|_{\rm HS}^2\). These are the
inherited Hessian identities, not a Sobolev embedding. One first
interchanges Hessian slots and then commutes on \(\nabla v\); no
derivative of curvature is introduced.

Equations (7), (10) give \(\int f_j^2\le C_L V(1+\kappa)+o(1)\). The
determinant Chern--Weil identity and (9) imply \[
\|F_A\|_2^2=2\|F_A^+\|_2^2-4\pi^2a^2
\le4\int f_j^2+4\pi^2(-a^2)_+.
\] The curvature of the positive-spinor connection is bounded in norm by
the \(W^+,s,\operatorname{Ric}^{\circ},F_A\) blocks; \(W^-\) acts
trivially on this representation. Its squared integral is therefore at
most \(C_LV(1+\kappa+b_a)+o(1)\). The spinor equation and \(|\Phi|\le2\)
likewise bound \(\int|\nabla_A^*\nabla_A\Phi|^2\) by
\(C_LV(1+\kappa)+o(1)\).

Apply (18) to \(\Phi\), and twice differentiate the parallel quadratic
map \(\Phi\mapsto\omega\). The resulting terms are bounded by
\(C(|\Phi||\nabla_A^2\Phi|+|\nabla_A\Phi|^2)\). Finally use
\(\Delta_H\omega=\nabla^*\nabla\omega+(s/3)\omega-2W^+\omega\). This
proves \[
B_\omega:=\fint|\Delta_H\omega|^2
\le C_L(1+\kappa+b_a)+o(1),
\] \[
\boxed{(1+B_\omega)d_\beta
\le C_L\kappa(1+\kappa+b_a)+o(1).}
\tag{19}
\] All the forms are smooth actual solutions. None of these estimates
differentiates the smoothing weight.

\hypertarget{spectral-compatibility-and-completion-of-the-proof}{%
\subsection{6. Spectral compatibility and completion of the
proof}\label{spectral-compatibility-and-completion-of-the-proof}}

The bounded self-dual-form compatibility theorem states that, for
\(\|\omega\|_\infty\le K\) and bounded constant \(t\), \[
|t|\,|1-2t|
\le C\sqrt{d_t(1+B_\omega+d_t)},
\tag{20}
\] where \(d_t\) is as in (17). Its proof is recorded in full in
\href{round21_analytic.md}{the analytic criterion}, Sections 2, 3 and 6.
The mechanism relevant to its use here is worth making explicit. The
same Hessian estimate bounds second derivatives of \(A_{\omega,t}\). The
weak contracted Bianchi identity, tested against
\(J_t=\operatorname{div}W(A_{\omega,t})\), gives \[
\fint|J_t|^2\le C|t|\sqrt{d_t(1+B_\omega+d_t)}.
\] If \(v=|\omega|^2\), the polynomial pointwise contraction gives \[
t^2v(|\nabla\omega|^2+|d\omega|^2)
\le |J_t|^2+\tfrac43t^2|dv|^2,
\quad
\fint|dv|^2\le C\sqrt{d_t(1+B_\omega+d_t)}.
\] Testing the self-dual Hodge equation against \(v\omega\) gives an
algebraic term \(16t^2(1-2t)/3\); all its derivative terms are
controlled by the preceding two displays. The zeroth-order errors cost
at most \(Ct^2\sqrt{d_t}\). Thus \[
\left|\frac{16}{3}t^2(1-2t)\right|
\le7\fint|J_t|^2+C t^2\sqrt{d_t(1+B_\omega+d_t)}.
\] Division by \(|t|\), with the case \(t=0\) separate, proves (20). The
argument neither differentiates the Ricci error nor assumes
\(d\omega=0\).

Substitute the produced form, with \(t=1/\beta\), into (20). Using
(17)--(19), then letting its smoothing allowance tend to zero, \[
\frac{(\beta-2)^2}{\beta^4}
\le C_L\kappa(1+\kappa+b_a).
\tag{21}
\] Absorb \(\beta^4\le L^4\) into the constant. When \(\kappa\le1\), the
right side is bounded by \(2C_L(1+b_a)\kappa\). When \(\kappa>1\), (2)
is immediate after choosing \(c_L\le(L-2)^{-2}\). This proves (2).

For (3), Chern--Weil gives \[
4\pi^2T=\int\left(2|W^+|^2+s^2/24
-|\operatorname{Ric}^{\circ}|^2/2\right)
\ge-\tfrac12 Vr.
\] Hence \(a^2\ge T\) implies \(b_a\le r/(8\pi^2)\). For \(\kappa\le1\)
this is bounded by \(1/(8\pi^2)\); for \(\kappa>1\) use the same trivial
split. Equation (4) is (3) with \(r=0\), after multiplying by
\(\beta V/(\beta+1)\).

\hypertarget{the-allowed-endpoint-without-a-scalar-curvature-band}{%
\subsection{7. The allowed endpoint without a scalar-curvature
band}\label{the-allowed-endpoint-without-a-scalar-curvature-band}}

The same construction has a structural consequence even when the
parameter approaches the allowed endpoint. It is not obtained by
declaring that endpoint forbidden.

Let \(H\) be the harmonic self-dual representative of \(a_g^+\), with
convention \(\operatorname{proj}_{\mathcal H^+}(iF_A^+)=2\pi H\). Set \[
\zeta_\beta=-\frac{3\pi\beta}{\beta+1}H.
\] \textbf{Theorem 2 (actual harmonic alignment).} Under the hypotheses
of Theorem 1, if \(\kappa\le1\), the same smooth forms can be chosen
with arbitrarily small allowance \(\eta>0\) so that \[
\fint|\omega-\zeta_\beta|^2\le C_L\kappa+\eta,
\quad
\fint|\nabla\omega|^2
\le C_L\sqrt{\kappa(1+\kappa+b_a)}+\eta.
\tag{22}
\] The amplitude, polynomial Weyl fit and Hodge-Laplacian bounds remain
(15), (17) and (19). In particular, at \(\beta=2\), the harmonic target
is \(-2\pi H\), and \(\kappa=r+144\pi^2D_{\rm LB}/V\). This conclusion
needs neither a scalar band nor a scalar spectral gap.

\textbf{Proof.} Retain the actual curvature moment from Section 3. Since
\(C\le E_Q\), its Hodge remainder satisfies \[
18\|iF_A^+-2\pi H\|_2^2
=C_{\rm act}-72\pi^2p\le S+o(1).
\] The curvature equation therefore gives \[
\|f_j\psi+4\sqrt2\pi H\|_2^2\le4S/9+o(1).
\tag{23}
\] Furthermore \[
\int(f_j-\sqrt{q_\beta})^2
\le2\int(f_j-\sqrt Q)^2
+2\int(\sqrt Q-\sqrt{q_\beta})^2
\le C_LV\kappa+o(1).
\] Here \((f_0-\sqrt Q)^2\le f_0^2-Q\le9e^2/16\), and the square-root
map is uniformly Lipschitz on \([3,\infty)\). Use \(|\psi|\le4\) and
multiply (23) by \(c_\beta^2/q_\beta\). The identity
\(4\sqrt2\pi c_\beta/\sqrt{q_\beta} =3\pi\beta/(\beta+1)\) proves the
first estimate in (22).

For the second, the derivative conclusion of the bounded-form theorem on
\(1/L\le t\le1/2\) is
\(\fint|\nabla\omega|^2\le C_L\sqrt{d_t(1+B_\omega+d_t)}\) when
\(d_t\le1\). If \(\kappa\) is small, (17) ensures this. For the
remaining compact range of \(\kappa\le1\), the integrated Hodge equation
and the bound on \(B_\omega\) give
\(\fint|\nabla\omega|^2\le C_L\sqrt{1+b_a}\); enlarging the constant
recovers (22). All allowances are taken small simultaneously for the
same solution.

The harmonic form need not have a uniform maximum bound. Thus (22) does
not license replacing \(\omega\) by \(\zeta_\beta\) inside the quadratic
Weyl model. It does give a precise global spectral conclusion: if
\(\Pi_{>\lambda}\) is the spectral projection of the rough Laplacian on
self-dual forms, then for every \(\lambda>0\), \[
\fint|\Pi_{>\lambda}\zeta_\beta|^2
\le C_L\left[\kappa+
\frac{\sqrt{\kappa(1+\kappa+b_a)}}{\lambda}\right].
\tag{24}
\] Indeed, project \(\zeta_\beta-\omega\) and use
\(\|\Pi_{>\lambda}\omega\|_2^2 \le\lambda^{-1}\|\nabla\omega\|_2^2\),
then let the allowance tend to zero. This controls an actual specified
harmonic class; it does not assert global Kähler topology or exact
metric correction.

At \(\beta=2\), there is a stronger direct estimate, independent of
\(b_a\). In the notation of (11), the retained Hodge identity is \[
\int|\nabla\psi|^2
=2\int|d\psi|^2+2\int E_\lambda
-\tfrac13\int(s+2w)\rho^2.
\] Since \(\rho\le4\), \(\int(s+2w)^2\le2\mathfrak D_2\), and
\(c_2^2=3/4\), this gives simultaneously \[
\fint|\nabla\omega|^2\le G(\kappa)+\eta,
\quad G(\kappa)=13\kappa+4\sqrt{2\kappa},
\quad
\fint|d\omega|^2\le13\kappa/2+\eta.
\tag{25}
\] The preceding harmonic-alignment estimate can likewise be made
explicit: \[
\fint|\omega+2\pi H|^2\le49\kappa+\eta.
\tag{26}
\] For (26), (6) gives \(\fint(q-6)^2\le18\kappa\). The cutoff estimate
implies \(\fint(f_j-\sqrt6)^2\le9\kappa+o(1)\), and (23) then costs at
most \(4\cdot9\kappa+S/(9V)\le49\kappa\). Thus in (24) the endpoint
right side can be replaced by \(98\kappa+2G(\kappa)/\lambda\), with no
square-cost hypothesis. The
\href{line1_composed_application.md}{endpoint application} gives this
statement and its spectral consequences separately.

\hypertarget{what-is-strengthened-and-what-is-not}{%
\subsection{8. What is strengthened, and what is
not}\label{what-is-strengthened-and-what-is-not}}

At \(\beta=2\), \[
\mathfrak D_2=144\pi^2D_{\rm LB},\qquad
D_{\rm LB}=\frac1{4\pi^2}\int(s^2/24+|W^+|^2/2)-p/2.
\] The classical equality set is negative cscK, compatible with the
selected Spin\(^c\) structure, not necessarily Einstein. Equation (2)
does not turn its zero deficit into an unconditional Ricci estimate.

At \(\beta=8\), Chern--Weil gives \[
\mathfrak D_8
=108\pi^2(T-2p/3)+\frac{27}{2}\int|\operatorname{Ric}^{\circ}|^2.
\tag{27}
\] Thus the old topological calibration is one member of the family. For
other parameters, (2) controls a different scalar--Weyl functional, with
its own genuine angular calibration supplied in Section 4. This is a
functional extension of the method, not a second discovery of the old
\(1/3\) blow-up obstruction.

An exact negative Kähler--Einstein metric with canonical class has
\(s=-4,w=2,p=V/(2\pi^2)\). Direct substitution yields \[
\frac{\mathfrak D_\beta}{V}
=\frac{4(\beta-2)^2}{\beta}.
\tag{28}
\] The quadratic vanishing at the allowed endpoint is consequently the
correct order; the theorem does not claim an optimal coefficient.

For \(0<\beta\le2\), the classical scalar and scalar--Weyl bounds
already provide a stronger elementary bound. If \(x=\|s\|_2,y=\|w\|_2\),
then \(x\ge\sqrt{32\pi^2p}\) and \(x+y\ge\sqrt{72\pi^2p}\). Minimizing
\(x^2+\beta y^2\) on these constraints gives
\(\int(s^2+\beta w^2) \ge8\pi^2p(4+\beta)\). Thus this range does not
need the new compatibility argument. For \(\beta>2\), that scalar
constraint is inactive at the classical relaxed minimizer. The
additional obstruction in (2) is differential, through weak Bianchi and
the actual bounded form.

For the Yau26 negative input, \(r\le4\delta_-(g)^2\). Under
\(a^2\ge T\), (3) therefore gives the usable estimate \[
4\delta_-(g)^2+\mathfrak D_\beta/V
\ge c'_L(\beta-2)^2.
\tag{29}
\] To obtain a new purely topological application one must additionally
force \(\mathfrak D_\beta/V\) small from independently justified
topological or geometric data. No such new topology is asserted here.
Neither (29) nor the allowed-root form supply proves existence of an
Einstein metric on an arbitrary smooth manifold.

\hypertarget{sources-and-verification}{%
\subsection{Sources and verification}\label{sources-and-verification}}

\begin{itemize}
\tightlist
\item
  C. LeBrun, \emph{Ricci Curvature, Minimal Volumes, and Seiberg--Witten
  Theory}, \href{https://arxiv.org/abs/math/0003068}{Theorem 2.3 and
  Propositions 3.1, 4.1}: the classical curvature inequalities and cscK
  equality set.
\item
  C. LeBrun, \emph{Four-manifolds, curvature bounds, and convex
  geometry}, \href{https://arxiv.org/abs/math/0611450}{Lemma 3.7}:
  weighted equations.
\item
  \href{round4_rigidity.md}{Integral truncated-weight proof}, Sections
  2--4: the inherited cutoff and moment estimates, reproduced above.
\item
  \href{round21_analytic.md}{Bounded-form spectral theorem}: the
  analytic compatibility theorem, including its polynomial contractions.
\item
  \href{line1_scalar_band_remainder.md}{Scalar-band supply}: the prior
  allowed cscK branch, which is not itself a forbidden-spectrum
  application.
\end{itemize}

Exact Einstein saturation forcing Kähler--Einstein structure already
appears in LeBrun's Corollary 2.5; it is not a new classification here.
The quantitative persistence under integral Ricci error, and the actual
all-scalar form-production estimates, are the assertions at issue.

The parameterized all-scalar functional estimate (2) is the claim proved
in this note. Its finite algebra is checked separately; those checks do
not replace the gauge-solvability hypothesis or the analytic proofs. No
comprehensive priority claim over later literature is made.

\end{document}
